% =====================================================================
%  EXPERIMENTS
%  Revised version. Every number below was checked against the code /
%  data files that produced the results:
%    topologies/{abilene,geant,germany50}_sndlib.json
%    marl_routing/real_traffic.py  (temporal 70/30 split)
%    train_all_monaco.sh, eval_ns3_all_monaco.sh  (launchers)
%    evaluate_ns3.py, run_ns3_phase2.py, aggregate_ns3_util.py
%    ns-3-dev/scratch/abilene-validate/abilene-validate.cc
%  Items still needing YOUR numbers are marked  % >>> TODO
% =====================================================================
\section{Experiments}

This section defines the evaluation protocol used to compare the proposed
decentralized MARL routing policy against OSPF and against a centralized
single-agent GNN reference. The protocol is stated in full before any result is
reported, so that every comparison in Section~\ref{sec:results} can be traced to
a fixed set of held-out traffic matrices, a fixed baseline, and a fixed
statistical test.

\subsection{Computing Environment and Reproducibility}
\label{sec:setup-compute}

All experiments were run on a single Linux server (Rocky~9) with four Intel Xeon
E5-4620~v4 CPUs at $2.10$\,GHz ($40$ cores) and $512$\,GB of RAM. No GPU was
used: the policy networks are small and both the analytical surrogate and the
packet-level simulator are CPU-bound, so training and evaluation were
parallelised across cores instead ($18$ concurrent jobs, each capped to two
threads via \texttt{OMP\_NUM\_THREADS}). The software stack was Python~3.9,
ns-3.42, PyTorch~2.6.0, PyTorch~Geometric~2.6.1, NetworkX, and
Stable-Baselines3~2.7.1 (used only for the centralized PPO reference; the
decentralized MAPPO learner is our own implementation).

Training and evaluation are driven by two launcher scripts
(\texttt{train\_all\_monaco.sh}, \texttt{eval\_ns3\_all\_monaco.sh}) that fix
the full configuration of every run and write one log file per run, so each of
the $18$ trained models and each evaluation directory is traceable to the exact
command line, seed, and hyperparameters that produced it. The complete
configuration is reproduced in Tables~\ref{tab:setup} and~\ref{tab:hyper}; the
code state used for all reported results is commit
\texttt{3064d76}.  % >>> TODO: update to the final commit hash before submission

\subsection{Network Instances and Traffic}
\label{sec:setup-instances}

Table~\ref{tab:setup} records the three SNDlib instances and the per-network
quantities of the protocol. Topology, capacities, and propagation delays are
taken from the instance files unchanged. Demands are the measured directed time
series shipped with SNDlib: Abilene, $5$-minute matrices over six months;
GEANT, $15$-minute over four months; Germany50, $5$-minute over one day.

\begin{table}[h]
\centering
\small
\begin{tabular}{lccc}
\hline
& Abilene & GEANT & Germany50 \\
\hline
Nodes $N$ & $12$ & $22$ & $50$ \\
Links $|\mathcal{E}|$ (undirected) & $15$ & $36$ & $88$ \\
Link capacities (Gb/s) & $9.92\ (\times 14),\ 2.48\ (\times 1)$ & $40$ & $40$ \\
Propagation delay range (ms) & $0.66$--$10.96$ & $0.58$--$33.98$ & $0.13$--$1.26$ \\
Demand pairs $|\mathcal{D}|$ & $132$ & $462$ & $2450$ \\
Pairs carried in ns-3 & all & all & top $200$ \\
Demand granularity & $5$\,min & $15$\,min & $5$\,min \\
\hline
Training scales $\mathcal{Z}$ & $\{8,12,16\}$ & $\{3,5,7\}$ & $\{35,50,65\}$ \\
Evaluation scale & $12$ & $5$ & $35$ \\
Delay penalty $\beta$ & $0.5$ & $0.5$ & $0.1$ \\
Hop cap $\sigma_{\max}$ & uncapped & uncapped & $4$ \\
Candidate matrices per scale & $20$ & $20$ & $20$ \\
Test matrices retained $|\mathcal{M}|$ & $5$ & $5$ & $5$ \\
\hline
\end{tabular}
\caption{Network instances and per-network protocol settings. Capacities and
delays are read directly from the SNDlib instance files. The demand scales
differ by network because the measured demands are light relative to these
over-provisioned backbones (Section~\ref{sec:setup-scaling}).}
\label{tab:setup}
\end{table}

\paragraph{Demand scaling.}
\label{sec:setup-scaling}
Measured demand on all three backbones is light relative to installed capacity,
so shortest-path routing is trivially optimal at the native magnitude and no
routing method can distinguish itself. Following standard practice in traffic
engineering studies, we retain the measured spatial and temporal \emph{structure}
of the matrices and scale only their overall \emph{magnitude} by a factor
$z\in\mathcal{Z}$ to set the operating point. Nothing in the demand pattern is
synthetic; $z$ is a single scalar applied uniformly to every entry of a measured
matrix.

\paragraph{Train/test separation.}
The demand time series of each network is split \emph{temporally}: the first
$70\%$ of the timeline forms $\mathcal{T}_{\mathrm{train}}$ and the last $30\%$
forms $\mathcal{T}_{\mathrm{test}}$. Within a split, candidate timestamps are
sampled evenly spaced across the interval. A temporal rather than random split is
used deliberately: randomly interleaved matrices five minutes apart are near
duplicates, and a random split would leak train information into the test set.

We state one limitation of this design plainly. The evaluation scale in
Table~\ref{tab:setup} lies \emph{inside} the training set $\mathcal{Z}$ for every
network. What is held out is therefore the \emph{traffic matrix}, not the load
regime: we test generalisation to unseen measured demand, not extrapolation to
unseen load magnitudes. Cross-scale and cross-topology generalisation are
treated separately in Section~\ref{sec:results-transfer}.
% >>> TODO: keep this pointer only if you keep the zero-shot transfer result.

\subsection{Methods Compared}
\label{sec:setup-methods}

Four routing methods are evaluated on identical inputs.

\begin{enumerate}
  \item \textbf{OSPF} (deployed baseline). The shortest path between each pair
        under the standard inverse-capacity link metric, with deterministic
        tie-breaking. On GEANT and Germany50 all links are $40$\,Gb/s, so this
        coincides with minimum-hop routing; on Abilene the single $2.48$\,Gb/s
        link receives a proportionally higher weight, which is the fair
        comparison, since a minimum-hop baseline would be penalised by a
        configuration choice no operator would make.
        % >>> IMPORTANT: your draft said "minimum-hop". Confirm which metric the
        % >>> exported ospf_path actually uses (marl_routing/*_routing_env.py calls
        % >>> nx.shortest_path with no weight => minimum-hop). If it is minimum-hop,
        % >>> either change the code to weight="inv_capacity" and re-export, or
        % >>> replace this paragraph with the minimum-hop wording AND report the
        % >>> inverse-capacity variant on Abilene as a robustness check.
  \item \textbf{Greedy best-response} (achievable reference). Each demand is
        placed, in turn, on the admissible path that minimises the resulting
        maximum link utilisation, using the same admissible set~\eqref{eq:mask}
        as the learned policies. This is not a global optimum, but it is a
        strong myopic reference that bounds how much of the available gain a
        learned policy actually captures.
  \item \textbf{Single-agent GNN} (architecture ablation). A centralized policy
        with full-network observability, trained with PPO. Comparing against it
        isolates the cost of decentralisation: it shares the MARL policy's
        backbone and reward, and differs only in observability and in being a
        single decision-maker.
  \item \textbf{MARL} (proposed). The decentralized per-node policy trained with
        MAPPO.
\end{enumerate}

The controlled variable is the path, and only the path. All four methods are
installed in ns-3 through the identical static-route mechanism, over the same
topology, the same demand matrix, the same flow set, and the same simulator
configuration; the routing tables are the sole difference between runs. Where a
flow cap is applied (top $200$ demands by rate on Germany50, needed to keep the
packet-level simulation tractable), the cap is computed once per matrix and the
resulting flow set is used identically by every method, so no method is
evaluated on a different workload from another.

\subsection{Simulation Configuration}
\label{sec:setup-ns3}

Every reported quality-of-service number is measured in ns-3 at packet level,
not by the analytical model used to shape the training reward. Each link is a
\texttt{PointToPointNetDevice} whose transmit queue is the default DropTail
queue of $100$ packets; no active queue management, shaping, or congestion
control is configured, so reported loss is pure buffer overflow. Traffic is
one-way UDP: each demand is a constant-rate \texttt{OnOff} application with a
$1400$-byte payload, so that loss and delay are attributable to the routing
decision and not to a transport-layer feedback loop. Propagation delays are
installed at nanosecond resolution, which matters on Germany50 where several
links are below $1$\,ms and would otherwise truncate to zero.

Each simulation covers $8$\,s of simulated time with all flows active over the
window $[2\,\mathrm{s},\,9\,\mathrm{s}]$, leaving a $2$\,s start-up margin before
measurement-relevant traffic begins. Because the byte counters run over the full
$8$\,s window while flows are active for $7$\,s, all utilisation figures are
divided by the exact, uniform factor $7/8$ to recover true utilisation; the
correction is identical for every method and every run. Demand rates and link
capacities are both divided by a common scale factor of $20$, which leaves every
utilisation, loss, and delay ratio invariant while reducing packet count enough
to make the full $18$-model evaluation feasible in wall-clock time.
% >>> TODO: state total wall-clock cost, e.g. "the full evaluation is N ns-3 runs
% >>> and completes in about H hours at 10-way concurrency". You can read this off
% >>> the timestamps in logs/ .

\begin{table}[h]
\centering
\small
\begin{tabular}{lcc}
\hline
& MARL (MAPPO) & Single-agent GNN (PPO) \\
\hline
Training steps & $600{,}000$ & $500{,}000$ \\
Learning rate & $3\times10^{-4}$ & $3\times10^{-4}$ \\
Discount $\gamma$ & $0.99$ & $0.995$ \\
GAE $\lambda$ & $0.95$ & $0.95$ \\
Clip range & $0.2$ & $0.2$ \\
Epochs per update & $10$ & $10$ \\
Minibatch & $512$ & $256$ \\
Entropy coefficient & $0.01$ & $0.01$ \\
Value coefficient & $0.5$ & $0.5$ \\
Gradient-norm clip & $0.5$ & $0.5$ \\
Hidden width & $128$ & $128$ \\
Candidate paths per pair $k$ & $3$ & $3$ \\
\hline
\end{tabular}
\caption{Learner hyperparameters. Values shared between the two methods were
deliberately held equal so that the comparison in
Section~\ref{sec:setup-methods} isolates decentralisation rather than tuning.}
\label{tab:hyper}
\end{table}

The per-network settings that do differ ($\beta$ and $\sigma_{\max}$ in
Table~\ref{tab:setup}) were selected on training matrices only and then frozen:
they are identical across all seeds and all test matrices, and no test matrix
influenced their choice.
% >>> TODO: if beta/sigma_max were in fact chosen by looking at test results,
% >>> say so honestly instead, or re-select them on a validation slice of the
% >>> train split. An examiner will ask.

\subsection{Congestion Regimes}
\label{sec:setup-regimes}

Results are reported separately for two congestion regimes, because the
routing problem is qualitatively different in each: when demand fits, OSPF is
already close to optimal and the achievable gain is small; when it does not,
path choice determines how much traffic is lost.

Let $G=(V,E)$ be the network with arc capacities $c_e$ (Mb/s), and let
$\mathcal{D}=\{(s,d)\}$ be the set of source--destination pairs with demands
$r_{sd}$. A routing $\pi$ assigns each pair a path $P^{\pi}_{sd}$, giving arc
loads and \emph{offered} utilisations
\begin{equation}
  \ell_e(\pi) \;=\; \sum_{(s,d)\in\mathcal{D}} r_{sd}\,
  \mathbb{1}\!\left[e \in P^{\pi}_{sd}\right],
  \qquad
  u^{\mathrm{off}}_e(\pi) \;=\; 100\cdot\frac{\ell_e(\pi)}{c_e}\quad[\%],
  \label{eq:util}
\end{equation}
and the congestion metric is the maximum offered link utilisation
\begin{equation}
  U^{\mathrm{off}}(\pi) \;=\; \max_{e\in E} u^{\mathrm{off}}_e(\pi).
  \label{eq:maxutil}
\end{equation}
Quantity~\eqref{eq:maxutil} is \emph{uncapped}: it is computed from demands and
may exceed $100\%$, which is precisely what makes ``how far past capacity a
routing pushes the network'' measurable. It is distinct from the utilisation
measured inside ns-3, written $U^{\mathrm{ns\text{-}3}}$, which counts bytes
actually carried per device and therefore saturates at $100\%$ by construction:
offered load beyond capacity appears as packet loss, not as utilisation above
line rate. Both are reported, and they are never conflated.

Test matrices are stratified by the offered utilisation that \emph{OSPF} would
incur, $U^{\mathrm{off}}_{\mathrm{OSPF}}$:
\begin{equation}
  \mathrm{regime}(r) \;=\;
  \begin{cases}
    \textbf{overload}, & U^{\mathrm{off}}_{\mathrm{OSPF}}(r) \;\ge\; 100\%,\\[2pt]
    \textbf{feasible}, & U^{\mathrm{off}}_{\mathrm{OSPF}}(r) \;<\; 100\%.
  \end{cases}
  \label{eq:regime}
\end{equation}
The label is a property of the traffic matrix under the baseline alone. It is
assigned before any learned policy is evaluated and does not depend on the
learned routing, so the stratification cannot favour the proposed method. From
the $20$ candidate matrices per network we retain the three most congested
overloaded matrices and the two most stressed feasible matrices
($|\mathcal{M}|=5$), i.e.\ the hardest instance of each regime rather than an
average one.

\subsection{Performance Metrics}
\label{sec:setup-metrics}

Three metrics are reported, all measured in ns-3.

\begin{enumerate}
  \item \textbf{Maximum link utilisation} $U^{\mathrm{ns\text{-}3}}$: the largest
        per-link utilisation observed in the simulation, computed from bytes
        enqueued per device and corrected by the $7/8$ active-window factor of
        Section~\ref{sec:setup-ns3}. It is the standard traffic-engineering
        objective and summarises how close the network is driven to its
        bottleneck.
  \item \textbf{Packet loss}: the fraction of offered packets not delivered to
        their destination, aggregated over all flows and expressed as a
        percentage. With DropTail queues and no congestion control, this is
        exactly buffer overflow caused by the routing decision.
  \item \textbf{Mean end-to-end delay}: the mean one-way delay of
        \emph{delivered} packets, in milliseconds.
\end{enumerate}

Delay must be read jointly with loss and never in isolation. Dropped packets
contribute no delay sample, so a routing that loses more traffic can report a
lower mean delay purely as an artefact of survivorship. We therefore report
loss and delay side by side and treat a delay reduction as meaningful only when
loss is equal or lower.

\subsection{Statistical Protocol}
\label{sec:setup-stats}

Every configuration is trained from three independent seeds ($0,1,2$), giving
$3\ \text{networks}\times 2\ \text{learned methods}\times 3\ \text{seeds}=18$
models. Each model is evaluated on the same $|\mathcal{M}|=5$ held-out matrices,
so a network--method cell contains $3\times5=15$ measurements per metric.

Two sources of variation are distinguished. The \emph{seed} varies only policy
initialisation and rollout sampling; it does not change the topology, the
demands, or the baseline. The \emph{matrix} varies the workload and is shared
across methods. OSPF and greedy best-response are deterministic and
routing-seed-independent, and are therefore evaluated once per matrix and
pooled, not re-seeded.

Because all methods are evaluated on the same matrices, every comparison is
\textbf{paired}, and we exploit that. For each metric we report:
\begin{enumerate}
  \item the mean $\pm$ standard deviation over the three seeds, per network and
        regime, which conveys the run-to-run stability of the learned policies;
  \item the paired per-matrix difference $\Delta = m_{\mathrm{MARL}} -
        m_{\mathrm{OSPF}}$ over the $15$ (seed, matrix) observations, summarised
        by its median and a $95\%$ bias-corrected bootstrap confidence interval
        ($10^{4}$ resamples, resampled at the matrix level to respect the
        pairing);
  \item a two-sided Wilcoxon signed-rank test on the same paired differences,
        with the exact $p$-value and the matched-pairs rank-biserial effect
        size.
\end{enumerate}
A bootstrap interval is preferred to a $t$-interval because $15$ observations
of a bounded, skewed quantity such as loss are not plausibly normal, and a
rank-based test is preferred for the same reason.

We do not run a significance test across seeds: $n=3$ is far too small to
support one, and reporting a $p$-value over three points would be misleading.
Seed variation is reported descriptively as a spread, and a result is described
as robust only when its sign is consistent across all three seeds. Where a
confidence interval crosses zero, or where the sign is inconsistent across
seeds, the result is reported as inconclusive rather than as a win.
% >>> TODO: these tests are DESCRIBED but not yet COMPUTED. Add a script that
% >>> loads results/*/summary.json, forms the 15 paired differences per
% >>> (network, regime, metric), and emits median, 95% BCa CI, Wilcoxon p, and
% >>> rank-biserial r. Then cite the numbers in Section~\ref{sec:results}.
